\uuid{5OVG}
\exo7id{4542}
\titre{exo7 4542}
\auteur{quercia}
\organisation{exo7}
\datecreate{2010-03-14}
\isIndication{false}
\isCorrection{false}
\chapitre{Suite et série de fonctions}
\sousChapitre{Autre}
\cpge{mpsi-mp2i.derivabilite.a, pcsi.derivabilite.a, ptsi.derivabilite.a, bcpst1.derivation.a, mp-mpi.suites-series-fonctions.a, pc.suites-series-fonctions.a, psi.suites-series-fonctions.a}
\module{Analyse}
\niveau{L2}
\difficulte{}

\contenu{
\question{
Soit $f\in \mathcal{C}^{\infty}(\R)$. On définit la suite $(f_n)_{n\in\N^*}$ par
$f_n = {f}^{(n)}$ (dérivée $n$-ème).
On suppose que ${(f_n)}_{n\ge 1}$ converge uniformément vers $\varphi$.
Que peut-on dire de $\varphi$ ?
}
}
